\section*{Limitations}
While our Adaptive-RAG shows clear advantages in effectiveness and efficiency by determining the query complexity and then leveraging the most suitable approach for tackling it, it is important to recognize that there still exist potential avenues for improving the classifier from the perspectives of its training datasets and architecture. Specifically, as there are no available datasets for training the query-complexity classifier, we automatically create new data based on the model prediction outcomes and the inductive dataset biases. However, our labeling process is one specific instantiation of labeling the query complexity, and it may have the potential to label queries incorrectly despite its effectiveness. Therefore, future work may create new datasets that are annotated with a diverse range of query complexities, in addition to the labels of question-answer pairs. Also, as the performance gap between the ideal classifier in Table~\ref{tab:main} and the current classifier in Figure~\ref{fig:classifier} indicates, there is still room to improve the effectiveness of the classifier. In other words, our classifier design based on the smaller LM is the initial, simplest instantiation for classifying the query complexity, and based upon it, future work may improve the classifier architecture and its performance, which will positively contribute to the overall QA performance.
